\documentclass{IOS-Book-Article}

\usepackage{times}
\usepackage{soul}
\usepackage{graphicx}
\usepackage{amsmath,amssymb,amsfonts}
\usepackage{booktabs}
\usepackage{multirow}
\usepackage{tabularx}
\usepackage{microtype}

\newtheorem{definition}{Definition}
\newtheorem{example}{Example}
\newtheorem{theorem}{Theorem}

% -------------------------------------
% Document Start
% -------------------------------------
\begin{document}

\begin{frontmatter}

\title{Measuring the Downstream Decision Impact of Predicate Extraction Errors in Symbolic Legal Engines}
\runningtitle{Measuring the Downstream Decision Impact of Predicate Extraction Errors in Symbolic Legal Engines}

% Authors & Affiliations
\author[A]{\fnms{Linh} \snm{Pham Van}%
\thanks{Corresponding Author: Linh Pham Van, College of Engineering \& Computer Science, VinUniversity, Viet Nam; E-mail: linhpv.vinuni@gmail.com.}},
\author[A]{\fnms{Ha-Thanh} \snm{Nguyen}%
\thanks{Corresponding Author: Ha-Thanh Nguyen, College of Engineering \& Computer Science, VinUniversity, Viet Nam; E-mail: thanh.nh6@vinuni.edu.vn.}}
and
\author[B]{\fnms{Ken} \snm{Satoh}%
\thanks{Corresponding Author: Ken Satoh, Center for Juris-Informatics, ROIS-DS, Japan; E-mail: ksatoh@nii.ac.jp.}}
\runningauthor{L.V. Pham et al.}

\address[A]{College of Engineering \& Computer Science, VinUniversity, Viet Nam}
\address[B]{Center for Juris-Informatics, ROIS-DS, Japan}

% Abstract
\begin{abstract}
Neuro-symbolic legal reasoning pipelines couple an LLM, which extracts structured predicates from a fact pattern, with a symbolic rule engine that applies an auditable legal rule to reach a decision. Such pipelines are typically evaluated by the extractor's predicate-level accuracy or F1-score, which implicitly treats every predicate error as equally consequential. We show this assumption does not hold: because a symbolic rule tree encodes conditional structure, some predicate errors flip the final decision while others are absorbed as redundant, and standard extraction metrics cannot tell the two apart. Using the PYTHEN framework applied to hearsay classification on LegalBench, we introduce a predicate-level intervention procedure that corrects one erroneous predicate at a time, reruns the rule tree, and measures whether the decision changes, yielding a Decision Impact Rate (DIR) for each predicate, which we compare against standalone extraction F1 and use to refine an extraction-error taxonomy for this pipeline along a decision-relevance dimension. Across three LLMs, we find that extraction difficulty and decision impact diverge sharply: the hardest predicate to extract is also the most decision-critical, while a comparably hard predicate is almost always inconsequential, and nominally identical error types carry markedly different downstream risk across models. These results argue for rule-aware error attribution, rather than aggregate predicate accuracy alone, as a necessary complement to evaluating and verifying LLM-symbolic legal reasoning pipelines.
\end{abstract}

% Keywords
\begin{keyword}
Neuro-Symbolic Legal AI \sep Explainability and Verification \sep Error Attribution \sep Large Language Models \sep LegalBench \sep Hearsay Classification
\end{keyword}

\end{frontmatter}

% -------------------------------------
% Sections
% -------------------------------------
\section{Introduction}
\label{sec:intro}

Large Language Models (LLMs) are increasingly proposed as tools for legal reasoning tasks, from statutory question answering\cite{louis2024interpretable} to case outcome prediction\cite{jiang2023legal}. Yet the same properties that make LLMs attractive, which is its fluent, free-form generation over unstructured text, also make them poorly suited to the standards the legal domain demands\cite{mik2024caveatlectorlargelanguage}. A legal decision is not merely a label, it is expected to be traceable to the rule that produced it, defensible against a specific line of argument, and reproducible under scrutiny. Chain-of-thought prompting improves the appearance of reasoning\cite{yu2023exploring}, but the underlying inference remains implicit, and free-form justifications frequently contain unsupported logical leaps that are difficult to detect from the output text alone\cite{galitsky2026discourse}. For high-stakes legal applications, opacity of this kind is not a minor limitation but a barrier to adoption.

Neuro-symbolic architectures have been proposed as a response to this problem: an LLM is used to extract structured, human-interpretable facts (predicates) from a natural-language fact pattern, and a symbolic reasoner then applies an explicit, auditable rule structure to those facts to reach a decision\cite{yang2025neuro}. Because the rule structure is fixed and inspectable, the resulting decision path can, in principle, be verified independently of the model that produced it. In this paper, we instantiate this division of labour as PYTHEN, a neuro-symbolic pipeline for hearsay classification under the LegalBench benchmark\cite{guha2023legalbench}, and show that it is empirically beneficial: replacing an LLM's final-answer judgment with a symbolic rule tree driven by LLM-extracted predicates improves F1-score and reduces instability across model families, relative to both direct prompting and chain-of-thought prompting. An oracle experiment, in which gold predicates are fed directly into the rule tree, further shows that the symbolic component reaches perfect agreement with ground truth once its inputs are correct - indicating that essentially all of the residual error in the neuro-symbolic pipeline originates upstream, in the LLM's predicate extraction step.

This result, however, exposes a gap that standard evaluation practice does not close. Prior work in relation extraction and structured information extraction commonly evaluates systems using corpus-level aggregate measures, particularly precision, recall, and F1 variants\cite{olek2023evaluationf1scorerecent,chen2023unifiedviewevaluationmetrics}. Such metrics summarize performance over extracted structures or relation instances, rather than exposing quality differences for individual predicate slots. This assumption does not hold once predicates are consumed by a rule tree. A rule tree encodes conditional structure: some predicates sit on a decisive branch where a single incorrect value flips the final verdict, while others are redundant with correlated conditions, are short-circuited by an earlier branch, or belong to a path the specific fact pattern never reaches. An extraction error on the first kind of predicate is decision-critical; the same kind of error on the second kind of predicate is, for that instance, functionally inert. Two models with identical predicate-level F1-scores can therefore differ substantially in how many of their errors actually reach the final decision, and a model that appears weaker by aggregate F1 may in practice make fewer decision-altering mistakes. Aggregate predicate metrics cannot distinguish these cases, because they are computed independently of the rule structure the predicates feed into.

This paper addresses that gap by asking, for a neuro-symbolic legal reasoning pipeline, not only \emph{how often} an LLM extracts a predicate incorrectly, but \emph{whether that error matters}, that is, whether correcting it would change the system's output. We introduce a predicate-level intervention methodology: for each fact pattern, we take the LLM-extracted predicate set, correct one erroneous predicate at a time toward its gold value, and rerun the PYTHEN rule tree to observe whether the decision changes. This yields, for every predicate error observed in practice, a direct measurement of its downstream decision impact, which we compare against the predicate's standalone extraction F1-score. We further refine an extraction-error taxonomy for this pipeline by adding a rule-position dimension, distinguishing decision-critical extraction errors from decision-irrelevant ones, and use it to characterize where LLM extraction failures actually threaten the reliability of the final legal conclusion. The hearsay classification task from LegalBench\cite{guha2023legalbench}, with its validated PYTHEN rule tree, provides a controlled setting in which every predicate's position in the rule structure is known, making such an attribution analysis tractable.

The contribution of this paper is therefore a methodology and an empirical analysis for rule-aware error attribution in neuro-symbolic legal reasoning, rather than a new architecture or a new headline accuracy result. Concretely, we (i) propose a predicate-level intervention procedure for measuring the downstream decision impact of individual predicate extraction errors in a symbolic legal engine; (ii) quantify, for the hearsay classification task, the extent to which predicate-level F1 over- or under-states a predicate's practical importance to the final decision; and (iii) refine the existing extraction error taxonomy to incorporate this decision-impact dimension, offering a more actionable diagnostic for where LLM-symbolic pipelines should be hardened. In doing so, we argue that evaluating neuro-symbolic legal systems requires looking beyond predicate accuracy in isolation, toward how errors propagate through the specific rule structure that governs the decision.

The remainder of this paper is organized as follows. Section~\ref{sec:related} situates this work relative to symbolic and hybrid approaches to legal reasoning and to error analysis in neuro-symbolic systems. Section~\ref{sec:method} formalizes the PYTHEN rule tree, the predicate-level intervention procedure, and the refined error taxonomy. Section~\ref{sec:experiments} reports the resulting decision-impact measurements on the hearsay classification task and compares them against standalone predicate F1. Section~\ref{sec:discussion} discusses the implications for evaluating and hardening neuro-symbolic legal pipelines, and Section~\ref{sec:conclusion} concludes.

\section{Related Work}
\label{sec:related}

\subsection{Symbolic and Rule-Based Legal Reasoning}
\label{subsec:related-symbolic}
Formalizing statutory rules as executable logic has a long history in AI and law, beginning with Sergot et al.'s encoding of the British Nationality Act as a Prolog program \cite{sergot1986}, which showed that a hand-authored rule structure could reproduce legal determinations exactly once the relevant facts were supplied - the same separation of fact-finding from rule application that underlies the PYTHEN pipeline in this paper. PROLEG \cite{satoh2010proleg} develops this logic-programming tradition further: a Prolog-based reasoning system that formalizes the Japanese "presupposed ultimate fact" theory (\textit{Yoken-jijitsu-ron}) governing civil litigation, it represents rules as a positive-condition part together with explicit \textit{exception} predicates, so that it can still reach a defensible conclusion under Japanese civil code's burden-of-proof structure when some facts remain unknown, rather than requiring every condition to be resolved to true or false. This explicit treatment of exceptions and open facts goes beyond what a single, fixed Boolean rule tree such as PYTHEN's can express, and marks one concrete direction - alongside deeper rule hierarchies and sequential gating - in which the rule-aware error attribution methodology proposed in this paper could be extended to richer rule structures (Section~\ref{subsec:limitations}). Subsequent work extended this tradition to formal argumentation and defeasible reasoning more broadly, providing frameworks for handling conflicting rules, exceptions, and burdens of proof that a purely Boolean rule tree cannot express \cite{benchcapon2007}. Prakken and Sartor's survey of this line of work from an argumentation perspective \cite{prakken2020} notes that its central appeal, and also limitation, is that soundness and explainability depend on a correctly and completely specified logical formalization, an assumption that traditionally required both the rules and the facts to be entered by hand rather than extracted from unstructured natural-language text.

\subsection{Empirical Legal NLP and Large Language Models}
\label{subsec:related-llm}
The rise of LLMs has shifted attention toward automatically bridging this gap between free text and formal input. LegalBench \cite{guha2023legalbench} introduced a benchmark of 162 tasks - including the \texttt{hearsay} classification task used in this paper - constructed to measure legal reasoning capabilities directly relevant to practicing lawyers, and evaluated LLMs primarily via direct or chain-of-thought prompting for a final label. This evaluation paradigm inherits a known reliability problem: an analysis of chain-of-thought faithfulness \cite{turpin2023unfaithful} shows that such explanations can be systematically unfaithful to the process that actually produced a model's answer, so a fluent justification does not guarantee that the stated reasoning caused the stated conclusion. A related faithfulness study of LLMs used as solvers and autoformalizers in legal reasoning \cite{wang2026knowyourlimits} documents an analogous divergence, finding that benchmark accuracy and logical faithfulness diverge, and identifying failure patterns - such as reporting conclusions unsupported by the model's own formal derivation - that aggregate accuracy does not surface. This paper's central concern, that predicate-level F1 does not indicate decision-level reliability, is a structural analogue of this same divergence, applied to a neuro-symbolic pipeline rather than to free-form or autoformalized reasoning.

\subsection{Hybrid and Neuro-Symbolic Approaches}
\label{subsec:related-hybrid}
Closer to this paper's architecture, a growing body of work couples LLMs with an explicit symbolic or structured component rather than relying on the LLM's final-answer judgment alone. Text-guided legal knowledge graph reasoning \cite{li2021legalkg} casts legal provision prediction as reasoning over a knowledge graph guided by text representations, combining neural text understanding with structured graph inference. A multi-agent framework for verifiable legal reasoning \cite{sadowski2025verifiable} similarly separates knowledge extraction - formalizing statutes into verifiable intermediate representations - from a second phase that applies this formalized knowledge through symbolic inference, reporting substantial accuracy gains over direct prompting on tax-calculation tasks. This body of work establishes that routing an LLM's output through a symbolic or structured intermediate representation is broadly beneficial, consistent with what the neuro-symbolic pipeline in this paper also finds. However, these evaluations, like the aggregate LegalBench baselines above, measure only whether the final decision is correct, not which of the LLM's intermediate extraction errors were responsible when it is not. None of this prior work asks whether a given extraction error would have changed the final decision had it been corrected in isolation, which is the question this paper addresses through predicate-level intervention on a fixed, auditable rule tree.

\section{Methodology}
\label{sec:method}

This section formalizes the neuro-symbolic pipeline that couples an LLM predicate extractor with the PYTHEN rule tree, and introduces the predicate-level intervention procedure used to measure the downstream decision impact of individual predicate extraction errors. We first present the pipeline and the rule tree it relies on (Sections~\ref{subsec:task}-\ref{subsec:tree}). We then introduce the core methodological contribution of this paper: single-predicate intervention (Section~\ref{subsec:intervention}), the decision impact metrics derived from it (Section~\ref{subsec:metrics}), and a rule-aware refinement of an extraction-error taxonomy for this pipeline (Section~\ref{subsec:taxonomy}).

\subsection{Task and Dataset}
\label{subsec:task}
We use the \texttt{hearsay} subset of LegalBench, in which each instance is a natural-language fact pattern together with a binary gold label indicating whether the described evidence constitutes hearsay. The subset contains 94 fact patterns. Following Federal Rule of Evidence 801 (FRE 801), we decompose the hearsay determination into five binary predicates, which we manually annotated for every fact pattern to obtain a gold predicate table:
\begin{itemize}
    \item \textit{is\_oral\_assertion}: the evidence involves a spoken statement or oral assertion.
    \item \textit{is\_written\_assertion}: the evidence involves a written statement or document.
    \item \textit{is\_non\_verbal\_assertive\_conduct}: the evidence involves non-verbal conduct intended as an assertion.
    \item \textit{is\_out\_of\_court}: the statement or conduct was made outside the current trial or hearing.
    \item \textit{is\_introduced\_to\_prove\_truth}: the evidence is introduced to prove the truth of the matter it asserts.
\end{itemize}
With 94 fact patterns and 5 predicates each, the gold table contains \textbf{470 predicate slots} in total. This predicate-level annotation, which is independent of the final hearsay label, is what makes the intervention procedure in Section~\ref{subsec:intervention} possible: it allows any single predicate in a model's extraction to be corrected in isolation, holding all other extracted predicates fixed.

\subsection{Neuro-Symbolic Pipeline}
\label{subsec:pipeline}
The pipeline separates fact extraction from legal inference. Given a fact pattern $x$, an LLM extractor $\mathcal{E}$ maps $x$ to a predicate vector
\begin{equation}
\label{eq:extractor}
\hat{p} = \mathcal{E}(x) \in \{0,1\}^5,
\end{equation}
where each coordinate corresponds to one of the five predicates in Section~\ref{subsec:task}. The extractor is prompted to also emit, for every predicate it sets to true, a supporting text span from $x$. This reference is not used by the symbolic reasoner, but is used during error annotation to distinguish genuine extraction errors from spurious ones (Section~\ref{subsec:taxonomy}).

The predicate vector $\hat{p}$ is then passed, unchanged, to the PYTHEN rule tree $R$, a fixed, hand-authored symbolic reasoner that computes the final decision
\begin{equation}
\label{eq:rule-tree}
\hat{y} = R(\hat{p}) \in \{\text{Hearsay}, \text{Not Hearsay}\}
\end{equation}
together with an explicit trace of which nodes of the tree were visited and which predicates determined each branching decision. Because $R$ is fixed and independent of the LLM, the only component that varies across models in this pipeline is the extractor $\mathcal{E}$; any change in $R$'s output between models is fully attributable to differences in their extracted predicate vectors.

\subsection{Formalizing the PYTHEN Rule Tree}
\label{subsec:tree}
The PYTHEN rule tree for hearsay classification encodes the statutory structure of FRE~801 as a Boolean function over the five predicates. An item of evidence is hearsay if and only if: (i) it constitutes an assertion, in oral, written, or non-verbal assertive form; (ii) that assertion was made out of court; and (iii) it is offered to prove the truth of the matter asserted:
\begin{equation}
\label{eq:hearsay-rule}
R(p) = \big(p_{\text{oral}} \lor p_{\text{written}} \lor p_{\text{nonverbal}}\big) \land\, p_{\text{out-of-court}} \land\, p_{\text{truth}}.
\end{equation}
Unlike a flat logistic combination of predicates, this structure is a decision tree: the disjunctive assertion clause is evaluated first, and the two conjunctive conditions act as gates that can independently veto a positive assertion clause. This structure is what makes different predicates unequally important for a given fact pattern, which we formalize next.

\begin{definition}[Structural Criticality]
\label{def:structural}
For a fact pattern with gold predicate vector $p$, predicate $p_j$ is \emph{structurally critical} if flipping $p_j$ alone, with all other predicates held at their gold values, changes the value of $R(p)$. Predicate $p_j$ is \emph{structurally redundant} otherwise.
\end{definition}

Under Equation~\ref{eq:hearsay-rule}, for instance, $p_{\text{written}}$ is structurally redundant whenever $p_{\text{oral}} = 1$ or $p_{\text{nonverbal}} = 1$ already satisfies the disjunctive clause, since flipping $p_{\text{written}}$ cannot then change $R(p)$; but $p_{\text{out-of-court}}$ and $p_{\text{truth}}$ are always structurally critical when the assertion clause holds, since either one alone can veto the decision. Definition~\ref{def:structural} is a static, ground-truth-only property of the rule tree and the instance.
 
\subsection{Predicate-Level Intervention}
\label{subsec:intervention}
This section operationalizes the analogous, empirical question for an LLM's actual predicate errors.

Standard evaluation of the pipeline compares $\hat{y}_i = R(\hat{p}_i)$ against the gold label $y_i = R(p_i)$ for each instance $i$, and separately reports predicate-level accuracy of $\hat{p}_i$ against $p_i$. Neither measurement indicates which of the (possibly several) predicate errors in $\hat{p}_i$ was responsible for a wrong final decision. We address this with a single-predicate intervention procedure.

For each instance $i$ where the extractor makes at least one predicate error, let $E_i = \{j : \hat{p}_{i,j} \neq p_{i,j}\}$ be the set of erroneous predicate indices. For every $j \in E_i$, we construct an intervened predicate vector $\hat{p}_i^{(j)}$ by correcting only predicate $j$ to its gold value and leaving every other predicate - including any other errors - at its extracted value:
\begin{equation}
\label{eq:intervention}
\hat{p}_{i,k}^{(j)} =
\begin{cases}
p_{i,j} & \text{if } k = j \\
\hat{p}_{i,k} & \text{otherwise.}
\end{cases}
\end{equation}
We rerun the (unmodified) PYTHEN rule tree on the intervened vector to obtain $\hat{y}_i^{(j)} = R(\hat{p}_i^{(j)})$, and compare it to the pipeline's original decision $\hat{y}_i = R(\hat{p}_i)$. This isolates the causal effect of correcting exactly one predicate error, holding the extractor's other outputs - correct or incorrect - fixed, and requires no retraining or re-prompting of the LLM: only the symbolic stage is rerun.

\subsection{Decision Impact Metrics}
\label{subsec:metrics}
We define the \emph{decision impact} of a single predicate error as a binary indicator,
\begin{equation}
\label{eq:impact-indicator}
\mathrm{Impact}(i,j) = \mathbb{1}\big[\hat{y}_i^{(j)} \neq \hat{y}_i\big], \qquad j \in E_i,
\end{equation}
which is 1 exactly when correcting that single error would have flipped the pipeline's final decision, and 0 when the error was decision-irrelevant given the model's other outputs on that instance. Aggregating Equation~\ref{eq:impact-indicator} over all observed errors of a given predicate $j$ across the dataset yields the \emph{Decision Impact Rate} (DIR) of that predicate,
\begin{equation}
\label{eq:dir}
\mathrm{DIR}(j) = \frac{1}{|\{i : j \in E_i\}|} \sum_{i \,:\, j \in E_i} \mathrm{Impact}(i,j),
\end{equation}
the fraction of that predicate's extraction errors that actually alter the model's final hearsay decision. We report $\mathrm{DIR}(j)$ alongside the standalone extraction $F_1$-score of predicate $j$ (computed as in our earlier predicate-extraction evaluation) for each of the three LLMs, which lets us directly test whether a predicate's extraction difficulty, as measured by $F_1$, is a reliable proxy for its practical importance to the downstream legal decision, or whether the two diverge - e.g., a predicate with comparatively high $F_1$ but high DIR would be a higher-priority target for improvement than its $F_1$-score alone suggests, and vice versa.

We additionally compute an instance-level counterpart, the \emph{decision-flip rate}: the fraction of originally misclassified instances ($\hat{y}_i \neq y_i$) for which correcting \emph{all} erroneous predicates in $E_i$ restores $\hat{y}_i = y_i$ when they are corrected one at a time versus jointly, which lets us check for interaction effects between simultaneous predicate errors that single-predicate intervention alone cannot capture.

\subsection{Refined Error Taxonomy}
\label{subsec:taxonomy}
We categorize every extraction error using a three-part type taxonomy (E1: Predicate Hallucination, E2: Predicate Omission, E3: Invalid Reference), and extend it with the decision-impact dimension defined above. Every extraction error is now annotated with a two-part label $(\text{Type}, \text{Impact})$, where \textit{Type} $\in \{\text{E1}, \text{E2}, \text{E3}\}$ is assigned as before from the extracted JSON, the gold predicate table, and the supporting reference span, and \textit{Impact} $\in \{\text{Critical}, \text{Benign}\}$ is assigned from $\mathrm{Impact}(i,j)$ in Equation~\ref{eq:impact-indicator}. This yields six refined error categories (e.g., \textit{Critical Hallucination} vs.\ \textit{Benign Hallucination}), which we use in Section~\ref{sec:experiments} to determine, for each model, what fraction of its extraction errors are merely statistical noise in the predicate table versus errors that materially undermine the reliability of the final legal decision.

\section{Experiments and Evaluation}
\label{sec:experiments}

\subsection{Experimental Setup}
\label{subsec:setup}
We evaluate three LLMs as predicate extractors, which are DeepSeek V3.2, GPT-OSS 120B, and Gemma 4 31B, on the 94-instance \texttt{hearsay} subset of LegalBench described in Section~\ref{subsec:task}, against the manually annotated gold predicate table. Each model is evaluated under three pipeline configurations: \textbf{Naive}, where the LLM is prompted directly for a binary hearsay verdict; \textbf{Chain-of-Thought (CoT)}, where the LLM produces a free-form chain-of-thought before its verdict; and \textbf{Neuro-symbolic (Ours)}, where the LLM only extracts the five predicates of Section~\ref{subsec:task}, which are then passed to the PYTHEN rule tree (Equation~\ref{eq:hearsay-rule}) to obtain the final decision.

\subsection{Baseline Context}
\label{subsec:overall}
Table~\ref{tab:baseline-comparison} reports the full comparative evaluation of the Neuro-symbolic pipeline against Naive and CoT prompting across all three models and four metrics. Routing predicate extraction through PYTHEN improves F1-score over both baselines for all three models (e.g., from 0.725 under CoT to 0.813 for GPT OSS 120B), and narrows the spread of F1-scores across models relative to Naive prompting: the Neuro-symbolic pipeline is also the only configuration in which every model's accuracy exceeds 0.77. This aggregate gain is the starting point for the present paper, whose question is orthogonal to it: \emph{given} that the Neuro-symbolic pipeline reduces errors on average, which specific predicate extraction errors are responsible for the decisions it still gets wrong, and which are inconsequential? The remainder of this section answers that question.

\begin{table}[htbp]
\centering
\caption{Comparative evaluation of Naive, Chain-of-Thought, and Neuro-symbolic (PYTHEN) pipelines across models.}
\label{tab:baseline-comparison}
\begin{tabular}{llcccc}
\toprule
\textbf{Approach} & \textbf{Model} & \textbf{Accuracy} & \textbf{Precision} & \textbf{Recall} & \textbf{F1-Score} \\
\midrule
\multirow{3}{*}{Naive LLMs} & DeepSeek V3.2 & 0.691489 & 0.603448 & 0.853659 & 0.707071 \\
 & GPT OSS 120B & 0.648936 & 0.666667 & 0.390244 & 0.492308 \\
 & Gemma 4 31B & 0.734043 & 0.673913 & 0.756098 & 0.712644 \\
\midrule
\multirow{3}{*}{LLMs with Chain-of-Thought} & DeepSeek V3.2 & 0.744681 & 0.634921 & \textbf{0.97561} & 0.769231 \\
 & GPT OSS 120B & 0.765957 & 0.74359 & 0.707317 & 0.725 \\
 & Gemma 4 31B & 0.776596 & 0.666667 & \textbf{0.97561} & 0.792079 \\
\midrule
\multirow{3}{*}{\shortstack[l]{Neuro-symbolic\\(LLMs with PYTHEN) \textit{(Ours)}}} & DeepSeek V3.2 & \textbf{0.776596} & \textbf{0.678571} & 0.926829 & \textbf{0.783505} \\
 & GPT OSS 120B & \textbf{0.819149} & 0.74 & \textbf{0.902439} & \textbf{0.813187} \\
 & Gemma 4 31B & \textbf{0.840426} & \textbf{0.75} & 0.95122 & \textbf{0.83871} \\
\bottomrule
\end{tabular}
\end{table}

\subsection{Oracle Experiment: Isolating the Symbolic Component}
\label{subsec:oracle}
To determine whether the residual error of the Neuro-symbolic pipeline originates in the LLM extractor or in the rule tree itself, we feed the manually annotated \emph{gold} predicate table directly into PYTHEN, bypassing the LLM entirely. Accuracy, Precision, Recall, and F1-score are all $1.0$. Because the rule tree in Equation~\ref{eq:hearsay-rule} reproduces every gold label exactly when given correct predicates, we conclude that $R$ is logically sound with respect to this dataset, and that all remaining error of the Neuro-symbolic pipeline is attributable to the predicate extraction step $\mathcal{E}$, not to symbolic inference. This result licenses the predicate-level intervention analysis of Section~\ref{subsec:dir-results}: since the rule tree itself is error-free, any decision flip observed under intervention can be attributed entirely to the extraction error being corrected.

\subsection{Standalone Predicate Extraction Performance}
\label{subsec:extraction}
Table~\ref{tab:predicate-extraction} evaluates $\mathcal{E}$ directly against the gold predicate table, over all $94 \times 5 = 470$ predicate slots. Gemma 4 31B is the strongest extractor on every metric (F1 = 0.894), followed by GPT-OSS 120B and DeepSeek V3.2, which post nearly identical F1-scores (0.820 and 0.819) via different error profiles: GPT-OSS 120B trades more false positives for fewer false negatives, while DeepSeek V3.2 does the reverse.

\begin{table}[htbp]
\centering
\caption{Standalone predicate extraction performance (470 predicate slots).}
\label{tab:predicate-extraction}
\begin{tabular}{lccccc}
\toprule
\textbf{Model} & \textbf{Accuracy} & \textbf{Precision} & \textbf{Recall} & \textbf{F1-Score} & \textbf{Macro F1} \\
\midrule
DeepSeek V3.2 & 0.845 & 0.789 & 0.851 & 0.819 & 0.689 \\
GPT-OSS 120B   & 0.843 & 0.775 & 0.871 & 0.820 & 0.729 \\
Gemma 4 31B   & \textbf{0.911} & \textbf{0.873} & \textbf{0.918} & \textbf{0.894} & \textbf{0.768} \\
\bottomrule
\end{tabular}
\end{table}

The gap between micro-averaged F1 and Macro F1 in Table~\ref{tab:predicate-extraction} is itself informative: Gemma 4 31B falls from 0.894 to 0.768, and DeepSeek V3.2 from 0.819 to 0.689, once every predicate is weighted equally rather than by its frequency in the dataset. Table~\ref{tab:per-predicate-f1} breaks this down by predicate. \textit{is\_oral\_assertion} and \textit{is\_out\_of\_court} are extracted reliably by all three models (up to 0.957 and 0.976 F1 for Gemma 4 31B), whereas \textit{is\_written\_assertion} is extracted poorly across the board (0.286-0.429 F1). Manual inspection indicates that this predicate requires distinguishing a genuine written assertion from textual elements that merely appear in a fact pattern without being offered as a statement of affirmation - a distinction models frequently fail to draw - compounded by the predicate's low base rate in the subset, which makes its F1-score highly sensitive to a small number of errors.

\begin{table}[htbp]
\centering
\caption{Predicate-level F1-score by model.}
\label{tab:per-predicate-f1}
\begin{tabular}{lccc}
\toprule
\textbf{Predicate} & \textbf{DeepSeek V3.2} & \textbf{GPT-OSS 120B} & \textbf{Gemma 4 31B} \\
\midrule
is\_oral\_assertion                & 0.897 & 0.929 & \textbf{0.957} \\
is\_written\_assertion             & 0.286 & 0.429 & 0.375 \\
is\_non\_verbal\_assertive\_conduct & 0.593 & 0.667 & 0.667 \\
is\_out\_of\_court                 & 0.910 & 0.877 & \textbf{0.976} \\
is\_introduced\_to\_prove\_truth   & 0.760 & 0.744 & \textbf{0.864} \\
\bottomrule
\end{tabular}
\end{table}

The disparity between overall F1 and per-predicate F1 already suggests that not all predicates contribute equally to extraction difficulty; Section~\ref{subsec:dir-results} tests the further, distinct question of whether this extraction difficulty ranking coincides with the ranking of predicates by their actual effect on the final decision.

\subsection{Predicate-Level Intervention: Decision Impact}
\label{subsec:dir-results}
We now apply the intervention procedure of Section~\ref{subsec:intervention} to every predicate error produced by each model on the \texttt{hearsay} subset. For every instance $i$ and every erroneous predicate $j \in E_i$, we correct $\hat{p}_{i,j}$ to its gold value, rerun the rule tree, and record $\mathrm{Impact}(i,j)$ (Equation~\ref{eq:impact-indicator}). Table~\ref{tab:dir} reports, for each predicate and model, the resulting Decision Impact Rate (DIR) alongside the predicate's standalone extraction F1-score from Table~\ref{tab:per-predicate-f1}; the subscript on each DIR value is the number of erroneous slots $|\{i : j \in E_i\}|$ it is computed over.

\begin{table}[htbp]
\centering
\caption{Decision Impact Rate (DIR) versus extraction F1 by predicate and model. Subscripts denote the number of observed errors the DIR is computed over.}
\label{tab:dir}
\begin{tabular}{lcccccc}
\toprule
& \multicolumn{2}{c}{\textbf{DeepSeek V3.2}} & \multicolumn{2}{c}{\textbf{GPT-OSS 120B}} & \multicolumn{2}{c}{\textbf{Gemma 4 31B}} \\
\textbf{Predicate} & F1 & DIR & F1 & DIR & F1 & DIR \\
\midrule
is\_oral\_assertion                 & 0.897 & $0.182_{(11)}$ & 0.929 & $0.125_{(8)}$  & 0.957 & $0.200_{(5)}$  \\
is\_written\_assertion              & 0.286 & $\mathbf{0.700_{(10)}}$ & 0.429 & $\mathbf{0.875_{(8)}}$  & 0.375 & $\mathbf{0.900_{(10)}}$ \\
is\_non\_verbal\_assertive\_conduct & 0.593 & $0.091_{(11)}$ & 0.667 & $0.000_{(9)}$  & 0.667 & $0.000_{(8)}$  \\
is\_out\_of\_court                  & 0.910 & $0.167_{(12)}$ & 0.877 & $0.188_{(16)}$ & 0.976 & $1.000_{(3)}$  \\
is\_introduced\_to\_prove\_truth    & 0.760 & $0.621_{(29)}$ & 0.744 & $0.394_{(33)}$ & 0.864 & $0.875_{(16)}$ \\
\bottomrule
\end{tabular}
\end{table}

Table~\ref{tab:dir} confirms the structural prediction of Definition~\ref{def:structural} and shows that extraction F1 is a poor proxy for decision importance. \textit{is\_written\_assertion} is the hardest predicate to extract for every model (F1 between 0.286 and 0.429) and simultaneously has, by a wide margin, the highest DIR of any predicate for every model (0.700-0.900): nearly every time a model gets this predicate wrong, that specific error flips the final hearsay decision. \textit{is\_non\_verbal\_assertive\_conduct} has comparably low F1 (0.593-0.667) but the lowest DIR of all five predicates (0.000-0.091): its errors are almost always structurally redundant, since \textit{is\_oral\_assertion} or \textit{is\_written\_assertion} already satisfies the disjunctive assertion clause of Equation~\ref{eq:hearsay-rule} for the same instance. The two predicates with the highest extraction F1, \textit{is\_oral\_assertion} and \textit{is\_out\_of\_court}, illustrate that F1 and DIR can diverge in the opposite direction: \textit{is\_oral\_assertion}'s rare errors are mostly benign (DIR $\leq 0.200$) because they are usually redundant with another true assertion predicate, whereas \textit{is\_out\_of\_court} - a conjunctive gate with no disjunctive backup - carries a markedly higher DIR for GPT-OSS 120B and, on a small sample, for Gemma 4 31B (all 3 of its errors are decision-critical). \textit{is\_introduced\_to\_prove\_truth}, the other conjunctive gate, is both a comparatively frequent error source and consistently high-DIR (0.394-0.875) across all three models, making it - together with \textit{is\_written\_assertion} - the predicate most responsible for the pipeline's residual decision errors. Overall, the rank correlation between extraction difficulty ($1 - F_1$) and DIR is weak to negative for four of the five predicates, and strongly positive only for \textit{is\_written\_assertion}: a predicate that is both hard to extract and disproportionately consequential when extracted incorrectly, making it the clearest target for extractor improvement identified by this analysis.

\subsection{Refined Error Taxonomy}
\label{subsec:taxonomy-results}
Applying the impact label from Section~\ref{subsec:taxonomy} to every E1-E3 extraction error over all $470$ predicate slots, Table~\ref{tab:taxonomy} reports, for each error type and model, how many errors are decision-critical versus benign, and what fraction of that error type is critical.

\begin{table}[htbp]
\centering
\caption{Refined error taxonomy: critical vs.\ benign extraction errors by model, out of 470 predicate slots.}
\label{tab:taxonomy}
\begin{tabular}{llccc}
\toprule
\textbf{Error Type} & \textbf{Impact} & \textbf{DeepSeek V3.2} & \textbf{GPT-OSS 120B} & \textbf{Gemma 4 31B} \\
\midrule
\multirow{3}{*}{E1: Predicate Hallucination} & Critical & 25 (5.32\%) & 23 (4.89\%) & 25 (5.32\%) \\
                                              & Benign   & 17 (3.62\%) & 26 (5.53\%) & 1 (0.21\%) \\
                                              & \% Critical & 59.5\% & 46.9\% & \textbf{96.2\%} \\
\midrule
\multirow{3}{*}{E2: Predicate Omission}      & Critical & 4 (0.85\%) & 1 (0.21\%) & 2 (0.43\%) \\
                                              & Benign   & 25 (5.32\%) & 24 (5.11\%) & 14 (2.98\%) \\
                                              & \% Critical & 13.8\% & 4.0\% & 12.5\% \\
\midrule
\multirow{3}{*}{E3: Invalid Reference}       & Critical & 1 (0.21\%) & 0 (0.00\%) & 0 (0.00\%) \\
                                              & Benign   & 1 (0.21\%) & 0 (0.00\%) & 0 (0.00\%) \\
                                              & \% Critical & 50.0\% & - & - \\
\bottomrule
\end{tabular}
\end{table}

The un-refined error distribution - E1/E2/E3 counts before the impact label is applied - shows E1 (hallucination) to be the dominant extraction error for all three models (18.1\% of predicate slots for DeepSeek V3.2, 13.8\% for GPT-OSS 120B and Gemma 4 31B), with E2 (omission) substantially rarer (4.3\%, 4.3\%, and 2.1\%, respectively); this distribution also includes 2 E3 instances for DeepSeek V3.2, where a predicate was set to true with no supporting reference text at all. Table~\ref{tab:taxonomy} refines this picture along the dimension that matters for reliability: \emph{what fraction of each error type actually reaches the final decision}. For DeepSeek V3.2 and GPT-OSS 120B, E1 errors are critical only 59.5\% and 46.9\% of the time, respectively - meaning close to half (or more, for GPT-OSS 120B) of their hallucinated predicates are absorbed by the rule tree's disjunctive structure and never affect the output. Gemma 4 31B breaks this pattern sharply: although it hallucinates roughly as often in absolute terms (25 critical instances, comparable to the other two models), 96.2\% of its E1 errors are decision-critical, because Gemma rarely hallucinates a redundant assertion predicate when another is already true - its few hallucinations are concentrated on predicates that matter. Across all three models, E2 (omission) errors are critical far less often than E1 errors (4.0-13.8\%), indicating that when these extractors miss a predicate, the missed predicate is disproportionately likely to be one the decision did not depend on. This asymmetry - hallucination being substantially more decision-relevant than omission - would not be visible from the E1/E2 frequency counts alone, and is only revealed by cross-referencing error type with the intervention-based impact label.

\section{Discussion}
\label{sec:discussion}

The predicate-level intervention results of Section~\ref{sec:experiments} support the central claim of this paper: not all predicate extraction errors affect the final legal decision equally, and their impact depends systematically on where the predicate sits in the symbolic rule structure. This section interprets what that divergence means for evaluating and deploying neuro-symbolic legal reasoning pipelines, discusses model-specific reliability differences it exposes, and states the limitations of the present study.

\subsection{Extraction Difficulty Is Not Decision Importance}
\label{subsec:interpretation}
Standard practice evaluates an LLM predicate extractor by its F1-score against a gold predicate table, implicitly treating extraction difficulty as a proxy for how much a given predicate's errors matter\cite{tan2026diagnosing}. Table~\ref{tab:dir} shows this proxy fails in both directions. \textit{is\_written\_assertion} is simultaneously the hardest predicate to extract (F1 between 0.286 and 0.429) and the one whose errors are most likely to flip the final decision (DIR between 0.700 and 0.900): here, difficulty and importance coincide, and this predicate deserves priority for extractor improvement. \textit{is\_non\_verbal\_assertive\_conduct} has comparably low F1 (0.593 - 0.667) but near-zero DIR (0.000 - 0.091): its errors are difficult to extract but structurally redundant, because the disjunctive assertion clause of Equation~\ref{eq:hearsay-rule} is almost always already satisfied by \textit{is\_oral\_assertion} or \textit{is\_written\_assertion} on the same instance. Conversely, \textit{is\_out\_of\_court} and \textit{is\_introduced\_to\_prove\_truth} - the two conjunctive gates in the rule - are extracted comparatively well, yet their rare errors carry disproportionate weight (DIR up to 1.000 and 0.875, respectively), because a conjunctive condition has no disjunctive backup to absorb a mistake. A practitioner who allocates extractor-improvement effort purely by predicate F1 would over-invest in \textit{is\_non\_verbal\_assertive\_conduct} and under-invest in the two conjunctive gates - exactly the misallocation that rule-aware error attribution is designed to prevent.

\subsection{From Diagnosis to Verification: Rule-Aware Triage}
\label{subsec:triage}
This divergence has a direct practical use that connects back to the verification role PYTHEN is intended to play in the broader neuro-symbolic pipeline: DIR can prioritize which extracted predicates are worth verifying, either by a human reviewer or by a secondary model call, before a decision is trusted. Rather than auditing all five predicates uniformly, or auditing whichever predicates the extractor is least confident about, a rule-aware verification policy would flag only extracted predicates that are both erroneous-prone and high-DIR for that instance's rule path - concretely, \textit{is\_written\_assertion} and \textit{is\_introduced\_to\_prove\_truth} in the hearsay setting, and, whenever the assertion clause is otherwise unsatisfied, \textit{is\_out\_of\_court}. Because low-DIR predicates such as \textit{is\_non\_verbal\_assertive\_conduct} rarely change the outcome even when wrong, they can be deprioritized without materially increasing the risk of an incorrect final decision. This reframes predicate-level verification from an undifferentiated checklist into a targeted, rule-path-dependent audit, reducing verification cost without a proportional loss of reliability - one concrete way the symbolic layer can support explanation and verification of LLM output in high-stakes legal settings.

\subsection{Model-Specific Reliability Profiles}
\label{subsec:model-differences}
Table~\ref{tab:taxonomy} shows that the same nominal error type carries different practical risk across models. DeepSeek V3.2 and GPT-OSS 120B's hallucinated predicates (E1) are decision-critical only 59.5\% and 46.9\% of the time, respectively - meaning a large share of their false positives are structurally absorbed and never reach the output. Gemma 4 31B breaks this pattern: despite hallucinating about as often in absolute terms, 96.2\% of its E1 errors are decision-critical. This is not visible from aggregate accuracy or F1 alone, where Gemma 4 31B is simply the strongest model on every metric (Table~\ref{tab:predicate-extraction}); the taxonomy reveals that when Gemma 4 31B does err, it tends to err on predicates the decision actually depends on, rather than on redundant ones. This suggests that predicate-level F1 rankings and decision-level reliability rankings need not coincide, and that a model selection or deployment decision informed only by the former risks underestimating how consequential a given model's remaining errors are in practice.

\subsection{Limitations}
\label{subsec:limitations}
\textbf{Single task and rule structure.} All results in this paper come from the \texttt{hearsay} subset of LegalBench, whose rule tree is a single shallow disjunction gated by two conjunctive conditions (Equation~\ref{eq:hearsay-rule}). LegalBench and the broader legal NLP literature contain tasks with substantially different rule shapes - deeper hierarchies\cite{lan2026hscodecomp}, multiple sequential gates\cite{fan2026lexam}, exception and override clauses\cite{cooper2026defab}, or predicates with numeric thresholds\cite{sadowski2025verifiable} - that have not been evaluated here. The specific divergence pattern we observe (a hard-to-extract predicate that is also high-impact, alongside a hard-to-extract predicate that is low-impact) is a property of this particular rule structure, and it remains an open question whether rule-aware error attribution surfaces the same kind of pattern, a different one, or none at all, on rule trees of a different shape. We treat this paper's findings as a case study establishing the methodology and its value on one task, not as a generalized claim about neuro-symbolic legal pipelines at large.

\textbf{Small-sample DIR estimates.} DIR is computed only over the errors a model actually makes, and for some predicate-model pairs this is a small count (e.g., 3 errors for \textit{is\_out\_of\_court} under Gemma 4 31B, 5 for \textit{is\_oral\_assertion}). A DIR of 1.000 estimated from 3 errors is directionally informative but statistically fragile, and should be read alongside the error count reported in Table~\ref{tab:dir} rather than as a precise probability.

\textbf{Single-predicate intervention does not capture interaction effects.} By construction (Equation~\ref{eq:intervention}), our intervention corrects one erroneous predicate at a time while holding all others - including other errors - fixed. When an instance has two or more simultaneous predicate errors, individually correcting either one may fail to flip the decision even though jointly correcting both would; such instances are undercounted as "no critical predicate found" under our procedure, even though the errors are jointly responsible for the wrong decision. A full combinatorial intervention over all subsets of $E_i$ would resolve this at the cost of exponential evaluation cost in $|E_i|$, which we did not pursue given that multi-error instances are a minority in our data.

\textbf{Gold predicate annotation and rule soundness are validated only on this dataset.} The oracle experiment (Section~\ref{subsec:oracle}) shows that the PYTHEN rule tree reproduces every gold label when given correct predicates for the \texttt{hearsay} subset, and the gold predicate table itself was manually annotated by the authors. Both the rule tree's logical soundness and the predicate annotations are therefore validated only within the scope of this dataset, not against independent legal experts or held-out cases.

\subsection{Future Work}
\label{subsec:future-work}
The most direct extension is to repeat the predicate-level intervention procedure on additional LegalBench tasks with structurally different rule trees - particularly ones involving exception clauses or sequential gating - to test whether the extraction-difficulty/decision-impact divergence observed here is a general property of neuro-symbolic legal pipelines or specific to the hearsay rule's shape. A second direction is to use DIR directly as a training or data-collection signal: weighting fine-tuning examples, few-shot demonstrations, or active-learning acquisition by a predicate's DIR rather than its base rate or current F1 would target extractor improvement at the errors most likely to change outcomes. Third, the rule-aware triage strategy outlined in Section~\ref{subsec:triage} could be implemented and evaluated as an actual human-in-the-loop verification interface, measuring whether DIR-prioritized review recovers most of the reliability of full predicate verification at a fraction of the review cost. Finally, extending the intervention procedure from single-predicate to joint (multi-predicate) corrections would address the interaction-effect limitation above and provide a more complete account of which \emph{combinations} of errors are jointly responsible for incorrect legal decisions.

\section{Conclusion}
\label{sec:conclusion}

This paper asked a question that aggregate evaluation of neuro-symbolic legal reasoning pipelines leaves unanswered: when an LLM extractor feeding a symbolic rule engine makes a predicate error, does that specific error actually change the resulting legal decision? Building on a neuro-symbolic pipeline that couples LLM predicate extraction with the PYTHEN rule tree for hearsay classification, and on an oracle experiment confirming the rule tree is logically sound given correct predicates, we introduced a predicate-level intervention procedure: correcting one erroneous predicate at a time, rerunning the rule tree, and recording whether the decision flips. This yields a Decision Impact Rate (DIR) for every predicate, which we compared directly against standalone extraction F1, and used to refine an extraction-error taxonomy for this pipeline along a rule-aware, decision-relevance dimension.

The results confirm the paper's central hypothesis: extraction difficulty and decision importance are not the same thing, and the rule structure determines how far apart they can be. \textit{is\_written\_assertion} is both the hardest predicate to extract and the one whose errors are most likely to flip the final decision, while \textit{is\_non\_verbal\_assertive\_conduct} is comparably hard to extract but structurally redundant, with errors that rarely reach the output. The refined taxonomy further showed that nominally identical error types carry different practical risk across models - Gemma 4 31B's hallucinated predicates are decision-critical 96.2\% of the time, versus 46.9-59.5\% for the other two models - a distinction invisible to standard predicate-level F1 or accuracy.

Taken together, these findings support treating rule-aware error attribution, rather than aggregate predicate accuracy alone, as a necessary component of evaluating and verifying LLM-symbolic legal reasoning pipelines: knowing that an extractor is imperfect is far less actionable than knowing which of its imperfections the final decision actually depends on. We reported this as a case study on a single LegalBench task with a specific rule structure, and Section~\ref{sec:discussion} outlines how the methodology should be extended to rule trees of different shapes, to joint (multi-predicate) intervention, and to concrete verification tooling built on the Decision Impact Rate. We hope this diagnose-which-errors-matter perspective proves useful beyond the hearsay setting studied here, wherever an LLM's outputs are consumed by an auditable symbolic reasoner rather than trusted directly.


% -------------------------------------
% References
% -------------------------------------
\bibliographystyle{vancouver}
\bibliography{references}

\end{document}
